\documentclass[11pt]{article}
% Publication-quality resume template (LuaLaTeX + latexmk recommended)
% US Letter, single-column, restrained typography

\usepackage{ifxetex}
\usepackage{ifluatex}
\ifnum\ifxetex 1\fi\ifluatex 1\fi=0
  \usepackage[T1]{fontenc}
  \usepackage[utf8]{inputenc}
\else
  \usepackage{fontspec}
  \defaultfontfeatures{Ligatures=TeX,Scale=MatchLowercase}
  % Prefer TeX Gyre fonts when available; fall back to Latin Modern when not installed
  \IfFontExistsTF{TeX Gyre Termes}{\setmainfont{TeX Gyre Termes}}{\setmainfont{Latin Modern Roman}}
  \IfFontExistsTF{TeX Gyre Heros}{\setsansfont{TeX Gyre Heros}}{\setsansfont{Latin Modern Sans}}
  \IfFontExistsTF{TeX Gyre Cursor}{\setmonofont{TeX Gyre Cursor}}{\setmonofont{Latin Modern Mono}}
\fi

\usepackage{microtype}
\microtypesetup{protrusion=false}
\usepackage[letterpaper,margin=0.7in]{geometry}
\setcounter{secnumdepth}{0} % disable section numbering
\usepackage{parskip}
\usepackage[hidelinks]{hyperref}
\usepackage{xurl}
\urlstyle{same}
\usepackage{xcolor}
\usepackage{enumitem}
\setlist[itemize]{topsep=2pt, parsep=0pt, itemsep=2pt, leftmargin=*, label=\textbullet}
\usepackage{graphicx}
\usepackage{needspace}
\setlength{\parskip}{5pt}

% Centralized typographic sizes
\newcommand{\sectionfontsize}{12pt}
\renewcommand{\baselinestretch}{1} % explicit 13pt leading is set below

% Pandoc macro definitions
\providecommand{\tightlist}{\setlength{\itemsep}{0pt}\setlength{\parskip}{0pt}}

% Header layout commands
\newcommand{\ResumeName}[1]{\par\noindent{\sffamily\bfseries\fontsize{23pt}{27pt}\selectfont #1}\par}
\newcommand{\BeginResumeHeader}[2]{%
  \noindent\begin{minipage}[t]{0.08\linewidth}\vspace{0pt}%
  \href{#1}{\includegraphics[width=0.85\linewidth]{icons/#2}}%
  \end{minipage}\hfill\begin{minipage}[t]{0.90\linewidth}\vspace{0pt}%
}
\newcommand{\EndResumeHeader}{\end{minipage}\par}
\newcommand{\ContactLine}[1]{\par\noindent{\normalfont\fontsize{10pt}{12pt}\selectfont #1}\par}
\newcommand{\ProfileLink}[3]{%
  \href{#2}{\raisebox{-0.15\height}{\includegraphics[height=9pt]{icons/#1}}\,#3}%
}
\newcommand{\CompactResume}{%
  \fontsize{10.5pt}{12pt}\selectfont
  \vspace*{0pt}%
  \setlength{\parskip}{2pt}%
  \setlist[itemize]{topsep=1pt,parsep=0pt,itemsep=1pt}%
  \renewcommand{\section}[1]{\par\Needspace{3\baselineskip}\vspace{4pt}\noindent{\sffamily\bfseries\fontsize{11pt}{13pt}\selectfont ##1}\par}%
  \renewcommand{\subsection}[1]{\par\Needspace{3\baselineskip}\vspace{2pt}\noindent{\sffamily\bfseries\fontsize{10.5pt}{12pt}\selectfont ##1}\par}%
}

% Section heading (avoid sectsty dependency; use simple unnumbered section formatting)
\renewcommand{\section}[1]{\par\Needspace{4\baselineskip}\vspace{7pt}\noindent{\sffamily\bfseries\fontsize{\sectionfontsize}{14pt}\selectfont #1}\par\smallskip}
\renewcommand{\subsection}[1]{\par\Needspace{5\baselineskip}\vspace{4pt}\noindent{\sffamily\bfseries\fontsize{11pt}{13pt}\selectfont #1}\par\smallskip}

% Ensure hyphenation and ragged-right for better extraction
\raggedright

\begin{document}
\pagestyle{plain}

% PDF metadata (from Pandoc top-level title/author set by the build script)
\hypersetup{
  pdftitle={Marco Del Fava},
  pdfauthor={Marco Del Fava},
  pdfcreator={LuaLaTeX + latexmk}
}

% Set consistent body font size
\fontsize{11pt}{13pt}\selectfont

% Main body from pandoc
\CompactResume
\ResumeName{Marco Del Fava}
\ContactLine{Research Software \& Neural Data}
\ContactLine{M.S. Biomedical Engineering student | UC Davis | Human intracranial electrophysiology}
\ContactLine{mdelfava@ucdavis.edu | 707-637-6544 | Sacramento, CA}

\ProfileLink{github.png}{https://github.com/BarkBarkBarkBarkBarkBarkBark}{GitHub}
·
\ProfileLink{linkedin.png}{https://www.linkedin.com/in/marc-del-fava}{LinkedIn}
·
\href{https://barkbarkbarkbarkbarkbarkbark.github.io/resumes_and_letters/}{Portfolio}

\section{Engineering Summary}\label{engineering-summary}

Biomedical Engineering M.S. student and neuroscience researcher at UC
Davis. I build Python acquisition and data-integrity workflows and use
coding agents with profiling, debugging, and regression tests.

\section{Technical Skills}\label{technical-skills}

\textbf{Scientific Programming:} Python, NumPy, Pandas, scikit-learn,
pytest, Git, Linux, Docker; MATLAB

\textbf{Neural Data \& Acquisition:} Neuralynx, Intan RHD, DarkHorse
MED, Neo, SpikeInterface, Kilosort4, MountainSort5

\textbf{Engineering Methods:} Scientific data parsing, metadata
validation and repair, waveform analysis, experimental synchronization,
profiling, regression tests

\section{Research Experience}\label{research-experience}

\textbf{Graduate Student}, UC Davis Health --- Z-Lab (Dr.~Jie Zheng) ---
Jul 2025 - Present

\begin{itemize}
\tightlist
\item
  Work with human intracranial electrophysiology recordings, developing
  Python workflows to prepare, process, and analyze single-neuron data.
\end{itemize}

\section{Selected Engineering
Projects}\label{selected-engineering-projects}

\subsection{Neural Data Formats and
Integrity}\label{neural-data-formats-and-integrity}

\begin{itemize}
\item
  Adapted a Neo-backed MED reader to skip corrupt parity channels whose
  metadata caused standard readers to abort, allowing processing without
  modifying source recordings.
\item
  Evaluated signal fidelity across 20 million samples on one channel (32
  kHz; 625 s), comparing original CSC and compressed MED recordings.
  Measured a maximum amplitude difference of 0.5 µV, with exact
  agreement after whole-microvolt source rounding.
\end{itemize}

\subsection{Neural Acquisition and Experimental
Synchronization}\label{neural-acquisition-and-experimental-synchronization}

\begin{itemize}
\item
  Built Python workflows for processing Intan recordings through
  SpikeInterface, including probe mapping, channel checks, spike
  sorting, waveform extraction, and quality metrics.
\item
  Developed synchronization software to align PsychoPy behavioral events
  with RHX recordings, using matched event labels and affine clock
  fitting to characterize drift between recording systems.
\end{itemize}

\subsection{Debugging Agent-Generated
Software}\label{debugging-agent-generated-software}

\begin{itemize}
\item
  Investigated a performance regression in Hosaka, a command-line
  software platform, where agent-generated schema validation made
  listing seven records take approximately ten seconds.
\item
  Used execution timing logs to identify repeated validation as the
  bottleneck, then introduced content-based caching and regression tests
  to preserve checks for changed or invalid schemas.
\end{itemize}

\section{Additional Professional
Experience}\label{additional-professional-experience}

\textbf{IT Administrator}, First Step Communities --- Jul 2021 - Present

\begin{itemize}
\tightlist
\item
  Maintain network and IT infrastructure across multiple shelter sites,
  including routing/VLAN segmentation, Microsoft 365 administration, and
  troubleshooting.
\end{itemize}

\textbf{Analyst}, Eurofins Test America --- Oct 2019 - Jun 2021

\begin{itemize}
\tightlist
\item
  Prepared soil and water samples using solvent extraction in a
  high-throughput environmental testing laboratory, following
  established quality assurance procedures.
\end{itemize}

\section{Education}\label{education}

\textbf{M.S. Biomedical Engineering}, University of California, Davis
--- 2026--Present

\textbf{M.S. Psychology}, Arizona State University --- 2023

\textbf{B.S. Neuroscience, Physiology, and Behavior}, University of
California, Davis --- 2017

\end{document}
